\section{Experiments}
\label{sec:experiments}

This section describes how the four encoder families of Section~\ref{sec:methodology} are evaluated: the recognition metrics used, the structure of the experiment matrix, and the attention-quality protocol that supplements the recognition numbers.

\subsection{Recognition Metrics}
\label{sec:metrics_reco}

Recognition quality is reported as \emph{character error rate} (CER) and \emph{word error rate} (WER), the standard metrics for HTR~\cite{marti2002iam, retsinas2022htrbestpractices}. For a predicted transcription $\hat{y}$ and reference $y$, CER is defined as
\begin{equation}
\label{eq:cer}
\mathrm{CER}(\hat{y}, y) \;=\; \frac{S + D + I}{N},
\end{equation}
where $S$, $D$, and $I$ are the number of character substitutions, deletions, and insertions required to transform $\hat{y}$ into $y$ (computed via Levenshtein alignment), and $N = |y|$ is the number of characters in the reference. CER is computed per line and averaged over the split. WER is the same quantity computed at the word level after whitespace tokenization. Both numbers are reported at the epoch with the best validation CER, following Retsinas~\etal~\cite{retsinas2022htrbestpractices}.

CER and WER alone do not answer the register-token question. In Section~\ref{sec:results} they are used to establish (a) the near-flat CER-versus-$R$ curve and (b) the seed-noise floor against which any register-count effect must be judged; the interpretability question is evaluated separately with the metrics of Section~\ref{sec:attn_eval}.

\subsection{Experiment Matrix}
\label{sec:matrix}

The evaluation is organized into six experiment blocks that together cover the recognition, ablation, and generalization questions. Table~\ref{tab:matrix} summarizes the blocks.

\begin{table}[h]
\centering
\small
\caption{Structure of the evaluation. IAM and READ2016 blocks use $150$-epoch schedules with the same seed and identical evaluation protocol; the synthetic block was capped by cluster walltime and its numbers are reported as lower bounds on convergence.}
\label{tab:matrix}
\begin{tabular}{lp{9cm}}
\toprule
Block & Purpose \\
\midrule
Register sweep (IAM)         & CNN--RNN baseline and ViT-RGTS at $R \in \{0, 2, 4, 8, 16\}$, with and without SWA. Isolates the register-token effect on recognition. \\
Pretrained lanes (IAM)       & TorchVision ViT-B/16 and TrOCR under three fine-tuning strategies (uniform LR, two-group differential LR, layer-wise LR decay, gradual unfreezing). Isolates the fine-tuning strategy from the pretrained-representation quality. \\
Architecture ablations       & CNN stem, head type, RNN depth, embedding dim, encoder depth, augmentation profile. Identifies which design choices are load-bearing. \\
Training-strategy ablations  & Base LR, batch/accumulation, SWA start epoch, SWA LR. Three additional seeds on the same configuration to establish the reproducibility floor. \\
READ2016 generalization      & Register sweep repeated on historical German. Tests whether the IAM findings are dataset-specific. \\
Synthetic pretraining        & All four families trained on ${\sim}950$K synthetic lines and evaluated on IAM test. Quantifies the render-to-real domain gap. \\
\bottomrule
\end{tabular}
\end{table}

Every ablation follows a single-variable discipline: change one variable from the ViT-RGTS baseline ($R = 4$, no SWA), keep every other seed and configuration identical, and evaluate on the same IAM test split with the same greedy CTC decoder. Any observed effect smaller than roughly twice the reproducibility floor (Section~\ref{sec:res_seeds}) is treated as unresolvable at this experimental budget and reported as such rather than over-interpreted.

\subsection{Attention-Quality Evaluation}
\label{sec:attn_eval}

The register-token intervention is designed to change what attention maps look like more than what CER measures. Three complementary attention-side evaluations are used in Section~\ref{sec:results}, applied to a fixed sample of eight IAM test lines chosen for diverse length, writer style, and character content, and re-used across all register counts to keep visual comparisons fair.

The first is a \textbf{per-character attention overlay} in the style of Nikolaidou~\etal~\cite{nikolaidou2024beyondmem}. For each recognized character $c$ at CTC decoding column $t_c$, the attention row from the corresponding patch-token query in the final Transformer layer is reshaped, upsampled to the input resolution, and overlaid on the input image using the \emph{inferno} colormap at $\alpha = 0.5$. The CTC decoding column $t_c$ is defined as the timestep at which the greedy decoder emits character $c$, and is illustrated for one word in Section~\ref{sec:res_bm}.

The second is a \textbf{register-absorption analysis}. For each Transformer layer $\ell$ the fraction of total attention mass directed at register tokens versus patch tokens, and the L2 norm of each token position in the final-layer output, are computed. If registers are absorbing the sink pattern as Darcet~\etal~\cite{darcet2024registers} predict for classification ViTs, register tokens should carry more norm than patches at $R > 0$.

The third is a \textbf{CTC-column concentration} measurement. For each recognized character the distance in patch-column units between the attention peak and the character's CTC decoding column is measured, along with the fraction of characters whose peak sits within two columns of the CTC column. This is the closest single-number analogue to a ``localization accuracy'' in the CTC self-attention setting. Raw self-attention is contaminated by fixed high-norm sink columns; a Gaussian column-prior centered at the CTC column is applied uniformly across all register configurations to suppress those before the peak is measured. The magnitude of this correction is itself a diagnostic: if registers were fully absorbing the sinks, the correction would be unnecessary.
